\scanpage{166}
\mspage{3}
% archive sheet number: 165
Dans la suite, on se propose de donner des renseignements sur (1.1.5) lorsque $X$ est de type fini sur $K$, $K$ étant un ``corps local'' comme à l'accoutumée\unsure{...}

1.2. Lorsque le \deleted{\illegible{2}} pt géom. $\xi$ de $X_K$ est localisé en un point rationnel sur $X$, alors la suite exacte (1.1.2) \unsure{splitte|splite} canoniquement, et \deleted{\unsure{que}} on peut écrire $\pi_1(X,\xi)$ comme un produit $\frac12$ direct
\[
\pi_1(X,\xi)\simeq\pi_1(X_{\bar K},\xi).I,
\tag*{(1.2.1)}
\]
moyennant une action de $I$ sur $\pi_1(X_{\bar K},\xi)$ :
\[
\deleted{\unsure{I\mathrel{\text{---}}}}\quad I\mathrel{\text{---}}\operatorname{Aut}(\pi_1(X_{\bar K},\xi)).
\tag*{(1.2.\unsure{2|4})}
\]
\unsure{De même on aura|De même aussi on a}
\[
\pi_1(X,\xi)'\simeq\pi_1(X_{\bar K},\xi)^{\unsure{(p')|(p)}}.I,
\tag*{(1.2.\unsure{3|8})}
\]
moyennant un hom.
\[
\deleted{\unsure{\pi}}I\mathrel{\text{---}}\operatorname{Aut}(\pi_1^{(p')}(X_{\bar K},\xi)).
\tag*{(1.2.4)}
\]
composé de \unsure{1.1.7|1.1.2} et l'hom de fonctorialité
\[
\operatorname{Aut}\pi_1\mathrel{\text{---}}\operatorname{Aut}\pi_1^{(p')}.
\]

\deleted{\unsure{Lem}} \unsure{Explicitons|Expliciter} ici la suite exacte, valable pour tt gpe profini $\pi$
\[
1\mathrel{\text{---}}\pi/\unsure{z|Z}\mathrel{\text{---}}\operatorname{Aut}(\pi)\longrightarrow\operatorname{Aut\,ext}\pi\deleted{^{\unsure{(p')}}}\mathrel{\text{---}}1
\qquad \{Z=\text{centre de }\pi.
\tag*{(1.2.5)}
\]
Appliquée au cas de $\pi_1^{(p')}$, cette suite nous montre ceci :
